\documentclass[10pt,a4paper]{amsart}
\usepackage[margin=19mm]{geometry}
\usepackage{amsmath,amssymb,mathtools}
\usepackage[scr=rsfs,cal=euler]{mathalfa}
\usepackage{xfrac}
\usepackage[unicode,hidelinks,bookmarks=false]{hyperref}
\newcommand{\ie}{i.e.\ }
\newcommand{\cf}{cf.\ }
\newcommand{\resp}{respectively\ }
\newcommand{\Subsection}{\subsection}
\providecommand{\nobreakdash}{\nobreak\hbox{-}}
\setlength{\emergencystretch}{2em}
\multlinegap0pt
\usepackage{amssymb}
\usepackage{bm}
\usepackage{tikz}
\usepackage{dsfont}
\usepackage{enumerate}
\newtheorem{lemma}{Lemma}
\newcommand{\eps}{\varepsilon}
\renewcommand{\leq}{\leqslant}
\title{Brownian behaviour of the Riemann zeta function around the critical line}
\author{Louis Vassaux}
\date{Source: arXiv:2505.07352v2 (CC BY 4.0)}
\begin{document}
\maketitle

\noindent\emph{Source and licence.} Brownian behaviour of the Riemann zeta function around the critical line. Version arXiv:2505.07352v2 (CC BY 4.0), distributed by arXiv under the Creative Commons Attribution 4.0 International licence, http://creativecommons.org/licenses/by/4.0/, as recorded in the arXiv licence metadata for this exact version and shown on the abstract page https://arxiv.org/abs/2505.07352v2. Full licence notice: \texttt{licenses/arxiv-2505.07352-CC-BY-4.0.txt}.\par\smallskip

\noindent\textit{Editorial scope.} Counted unit: the article's lemma tech (source0.tex lines 1109--1115). The article's complete original proof of it is delivered with it (source0.tex lines 1173--1238, one \texttt{proof} environment of 66 lines, which the article places away from the statement). No mathematics is written, repaired or re-derived: the statement and the proof are the article's own lines, in the article's own notation, and no retained line is substituted or re-flowed.  No further source-local context is needed: every cross-reference of every retained line resolves inside the two delivered spans.  Nothing was omitted from any retained span, so the package carries no omission record.  No retained line contains a citation, so the wrapper carries no bibliography and no bibliography entry.  No retained line contains a figure environment, an image-inclusion command or drawing code, so no image is delivered and the package declares no asset dependency.  Editorial formatting only, in addition: an AMS \texttt{amsart} preamble that restates the article's own declarations and macros used by the retained lines, namely the environments \texttt{lemma} on separate counters, together with the standard packages those lines need.  The retained environments print in this document as Lemma 1 in order of appearance, whereas the article prints the same items with the numbering of its own full document.  The complete native source of the article as distributed in the arXiv e-print for this exact version is bundled unchanged as \texttt{sources/177950/source0.tex}.
\begin{lemma}
\label{tech}
    Let $0 \leq \alpha \leq \beta \leq 1$, be functions of $T$, and set  $\eta = (\log T)^{-\alpha}$, $\eta' = (\log T)^{-\beta}$. Then,
    \begin{equation}
            \sum_{p \leq x^3 \text{ prime}} \frac{1}{p^{1 + \eta}} - \frac{1}{p^{1 + \eta'}} \ll (\beta - \alpha)\log \log T.
    \end{equation}
\end{lemma}

\begin{proof}[Proof of Lemma \ref{tech}]
First, note that by taking the derivative of $\alpha \mapsto \dfrac{1}{(\log T)^\alpha}$, we obtain

\begin{equation}
    \label{iaf}
    \eta - \eta' \ll \eta (\beta - \alpha) \log \log T.
\end{equation}
We shall be using equation (\ref{iaf}) throughout the rest of this paper.
Let us rewrite:
\begin{equation}
        \sum_{p \leq x^3} \frac{1}{p^{1 + \eta}}  = \sum_{n \leq x^3} \frac{\pi(n) - \pi(n-1)}{n^{1+\eta}} 
 = \sum_{n \leq x^3} \pi(n) (\frac{1}{n^{1 + \eta}} - \frac{1}{(n+1)^{1 + \eta}}) + \frac{\pi(N)}{N^{1+\eta}}.
\end{equation}
where $N = \lfloor x^3 \rfloor$, and $\pi$ denotes the prime-counting function. We now wish to approximate this sum by its associated integral. Specifically, setting

\begin{equation}
    I_x(\eta) = \int_1^{x^3} \pi(u)(\frac{1}{u^{1+\eta}} - \frac{1}{(u+1)^{1 + \eta}})du,
\end{equation}
we will split up our problem:
\begin{equation}
    \sum_{p \leq x^3} \frac{1}{p^{1 + \eta}} - \sum_{p \leq x^3} \frac{1}{p^{1 + \eta'}} = \left(I_x(\eta')  - \sum_{p \leq x^3} \frac{1}{p^{1 + \eta'}} \right) - \left(I_x(\eta)  - \sum_{p \leq x^3} \frac{1}{p^{1 + \eta}} \right) + (I_x(\eta) - I_x(\eta')).
\end{equation}
Now,
\begin{equation}
    \begin{split}
        \dfrac{d}{d\eta} & \left(I_x(\eta)  - \sum_{p \leq x^3} \frac{1}{p^{1 + \eta}} \right) \\
        & = (1 + \eta) \left( \int_{1}^{x^3}\pi(u) (-\frac{\log u}{u^{1 + \eta}} + \frac{\log \lfloor u \rfloor}{\lfloor u\rfloor^{1 + \eta}} + \frac{\log (u+1)}{(u+1)^{1+\eta}} - \frac{\log \lfloor u+1 \rfloor}{\lfloor u+1\rfloor ^{1+\eta}}) du - \frac{\pi(N) \log N}{N^{1 + \eta}} \right) \\
        & \ll \int_{1}^{x^3}\pi(u) \frac{\log u}{u^{3+\eta}}du + \frac{\pi(N)\log N}{N^{1+\eta}} \ll 1
    \end{split}
\end{equation}
since $\pi(u) \sim \dfrac{u}{\log u}$. As a result,
\begin{equation}
    \left(I_x(\eta)  - \sum_{p \leq x^3} \frac{1}{p^{1 + \eta}} \right) - \left(I_x(\eta')  - \sum_{p \leq x^3} \frac{1}{p^{1 + \eta'}} \right) \ll (\eta - \eta') \ll (\beta - \alpha) \log \log T
\end{equation}
by (\ref{iaf}). As a result, we now just need to control $I_x(\eta) - I_x(\eta')$:
\begin{equation}
\begin{split}
    I_x(\eta) - I_x(\eta') & \ll \int_{2}^{x^3} \dfrac{u}{\log u}\left(\frac{1}{u^{1 + \eta}} - \frac{1}{(u+1)^{1 + \eta}} - \frac{1}{u^{1 + \eta'}} + \frac{1}{(u+1)^{1 + \eta'}} \right) du \\
    & = \int_2^{x^3} \left(\frac{\eta}{u^{1 + \eta} \log u} - \frac{\eta'}{u^{1 + \eta'}\log u} \right) du + \int_2^{x^3} \eps(u, \eta) - \eps(u, \eta')du
\end{split}
\end{equation}

with $\eps(u, \eta) = \dfrac{1}{u^{\eta} \log u} - \dfrac{u}{(u+1)^{1 + \eta} \log u} - \dfrac{\eta}{u^{1 + \eta} \log u}$. Since $\dfrac{d}{d \eta} \eps(u, \eta) \ll \dfrac{1}{u^{2 + \eta}}$,

\begin{equation}
    \int_2^{x^3} (\eps(u, \eta) - \eps(u, \eta'))du \ll \eta - \eta' \ll (\beta - \alpha) \log \log T.
\end{equation}
In order to bound the main integral, we may now set $v = \eta \log u$ (resp. $v = \eta' \log u$):
\begin{equation}
    \begin{split}
        \int_2^{x^3} & \left(\frac{\eta}{u^{1 + \eta} \log u} - \frac{\eta'}{u^{1 + \eta'}\log u} \right) du  = \eta \int_{\eta \log 2}^{3 \eta \log x} \frac{dv}{v e^v} - \eta' \int_{\eta' \log 2}^{3\eta' \log x} \frac{dv}{v e^v} \\
        & = (\eta - \eta') \int_{\eta \log 2}^{3 \eta' \log x} \dfrac{dv}{v e^v} + \eta \int_{3 \eta' \log x}^{3 \eta \log x}\dfrac{dv}{v e^v} - \eta' \int_{\eta' \log 2}^{\eta \log 2} \dfrac{dv}{v e^v} \\
        & = I_1 + I_2 - I_3.
    \end{split}
\end{equation}
It is quite clear that $I_2 \ll I_3$; meanwhile,
\begin{equation}
    I_3 = \eta' \int_{\eta' \log 2}^{\eta \log 2} \frac{1}{v}(1 + O(1))dv \ll \eta' \log \frac{\eta}{\eta'} + (\eta - \eta') \ll (\beta - \alpha) \log \log T.
\end{equation}
Finally,
\begin{equation}
        I_1  \ll (\eta - \eta') \left( \int_{\eta \log 2}^1 \dfrac{dv}{v} + \int_{1}^{3 \eta' \log x} e^{-v} dv \right) 
\ll \eta (1 + \log \eta)(\beta - \alpha) \log \log T \ll (\beta - \alpha)(\log \log T),
\end{equation}
which concludes the proof.
\end{proof}

\end{document}
